\section{Introduction}
\label{sec:intro}

A growing family of systems claims to generate scientific research
questions, hypotheses, or ideas
\citep{lu2024aiscientist,wang2024scimon,baek2024researchagent}. How do we
know whether any of them is good at it? Today, essentially every
evaluation falls into one of three patterns: an \emph{expert score} (a
panel rates novelty and significance on a Likert scale), an \emph{LLM
score} (a language model rates the same properties), or a \emph{case
study} (a handful of generated ideas is narrated persuasively). All
three share the same defect: they are subjective. Expert panels disagree
with each other and with themselves \citep{si2024llmideas}; LLM judges
inherit the biases of their training distribution and can be steered by
phrasing; case studies are selected by the authors. None of these
evaluations can be \emph{wrong} in a way that data could demonstrate,
which is another way of saying none of them measures anything.

Science itself offers a harder criterion. Research questions are bets
about where inquiry should go next, and the scientific community
eventually settles those bets: it invests observing time, funds
follow-ups, writes papers that answer some questions, poses others
independently, and refutes the premises of a few. This suggests a simple
standard:

\begin{quote}
\emph{A scientific question is valuable if future science actually
invests in answering it.}
\end{quote}

The standard becomes an evaluation protocol the moment we rewind the
clock. Fix a historical cutoff. Give a system only the literature
available before the cutoff. Freeze the questions it generates. Then let
the literature published \emph{after} the cutoff---which the system
never saw---grade the bet: Was the question answered? Substantially
advanced? Independently posed by working scientists? Ignored? Was its
underlying premise confirmed, or refuted? We call this procedure
\textbf{historical backtesting}, by analogy with the evaluation of
trading strategies on held-out past data \citep{bailey2014backtest}, and
with recent forecasting benchmarks that score models against events
occurring after training \citep{zou2022autocast}.

\begin{figure}[t]
\centering
\begin{tikzpicture}[
    node distance=0.55cm,
    stage/.style={draw, rounded corners=2pt, minimum width=6.2cm,
                  minimum height=0.72cm, align=center, font=\small},
    past/.style={stage, fill=blue!8},
    future/.style={stage, fill=orange!12},
    eval/.style={stage, fill=green!10},
    arr/.style={-{Stealth[length=2.2mm]}, thick}]
  \node[past]   (p1) {Past Literature ($\leq$ cutoff)};
  \node[past,   below=of p1] (p2) {Generate Scientific Questions};
  \node[past,   below=of p2] (p3) {Freeze Questions};
  \node[future, below=of p3] (f1) {Future Literature (after cutoff)};
  \node[future, below=of f1] (f2) {Retrieve Evidence};
  \node[eval,   below=of f2] (e1) {Historical Outcome Assessment};
  \node[eval,   below=of e1] (e2) {Benchmark Metrics};
  \draw[arr] (p1) -- (p2);
  \draw[arr] (p2) -- (p3);
  \draw[arr] (p3) -- (f1);
  \draw[arr] (f1) -- (f2);
  \draw[arr] (f2) -- (e1);
  \draw[arr] (e1) -- (e2);
  \node[right=0.5cm of p2, font=\footnotesize\itshape, text width=4.4cm, align=left]
    {any system: evidence graph,\\ prompted LLM, heuristic, human};
  \node[right=0.5cm of f1, font=\footnotesize\itshape, text width=4.4cm, align=left]
    {temporally isolated;\\ never seen at generation time};
  \node[right=0.5cm of e1, font=\footnotesize\itshape, text width=4.4cm, align=left]
    {constrained judge +\\ human adjudication};
\end{tikzpicture}
\caption{The historical backtesting protocol. No question generator,
evidence graph, or particular LLM appears in the loop: the protocol
evaluates \emph{frozen question lists}, whatever produced them.}
\label{fig:protocol}
\end{figure}

Crucially, the protocol contains no question generator
(Figure~\ref{fig:protocol}). It takes a frozen list of questions as
input and returns outcome labels and metrics as output. Any
system---an evidence-graph pipeline, a prompted LLM, a citation
heuristic, a human scientist---can be evaluated under identical
conditions. This is what makes it a benchmark rather than a validation
appendix for one particular architecture.

\paragraph{Contributions.} We make three claims, in decreasing order of
strength:

\begin{enumerate}
  \item \textbf{Protocol (strong).} We formalize historical backtesting
    as an evaluation protocol for scientific question discovery:
    temporal isolation rules, a question-freezing requirement, fixed
    future-evidence retrieval, a two-dimensional outcome taxonomy that
    separates a question's fate (\labelfmt{answered},
    \labelfmt{partially\_addressed}, \labelfmt{posed\_but\_open},
    \labelfmt{not\_addressed}) from its premise's fate
    (\labelfmt{supported}, \labelfmt{refuted}, \labelfmt{weakened},
    \labelfmt{still\_plausible}, \labelfmt{not\_applicable}), and
    metrics defined over those labels (Sections~\ref{sec:protocol}
    and~\ref{sec:metrics}).
  \item \textbf{Benchmark instance (medium).} We release a reproducible
    astronomy instance with temporally isolated past and future corpora
    (2{,}512 and 1{,}891 papers; cutoff 2020-12-31), frozen questions,
    released retrieval records, adjudicated labels, one-command metric
    computation with CI-enforced reproducibility, a submission format,
    and four reference baselines (Sections~\ref{sec:dataset}
    and~\ref{sec:baselines}).
  \item \textbf{Empirical pilot (cautious).} In an initial set of ten
    questions generated from pre-cutoff evidence only, all ten were
    substantively engaged by the 2021--2026 literature, including one
    whose underlying premise was subsequently refuted by three
    independent analyses (Section~\ref{sec:results}). With $n=10$
    questions from one system in one domain, we do \emph{not} claim
    that AI can reliably identify the most valuable future scientific
    questions; we claim the protocol can tell us, eventually, whether
    it can.
\end{enumerate}

What we believe is ultimately most useful here is not any single
result but the change of category: scientific question quality moves
from a matter of taste (``this question seems interesting'') to a
measured, comparable quantity (``under historical backtesting, 100\% of
this system's questions were engaged by later literature; 10\% led to a
premise refutation''). Every artifact needed to run the protocol---data,
code, labels, and checks---is public, and Astronomy~v1 is offered as the
first instance of the benchmark, not as its definition.
